\section{Empirical results and discussion}\label{sec:results}

The results proceed in three steps.
First, we ask when the conditional envelope classifies firm pairs informatively.
Second, we estimate whether article-embedding distance is associated with
return distance after symmetric firm effects.
Third, we assess whether the sign survives a later return window and
alternative representations.

\subsection{Conditional W$_2$ envelope}\label{sec:fixed-w2-test}

How tightly must the characteristic-to-exposure map be restricted for the
conditional covariance envelope to classify observed firm pairs?
The envelope classifies pairs informatively only under tight restrictions:
the zero-slack isometric case is decisive, but informativeness deteriorates
rapidly once even modest transmission slack is admitted.
The envelope exercise evaluates \Cref{eq:random-kappa-bracket} with the
common 128-article exact W$_2$ matrix and standardized total-return moments.
A covered dyad has nonnegative transport excess throughout the
scenario-implied latent distance interval; a violated dyad has negative
excess throughout; and an unresolved dyad has an interval that spans zero.
The zero-slack isometric cell $(L,\tau)=(1,0)$ covers
\ppnum[3]{confound-w2-l1-tau0-coverage} of observed dyads and violates the
restriction for \ppnum[3]{confound-w2-l1-tau0-violation}, leaving no
dyads unresolved.
Even the modest relaxation $(L,\tau)=(1,0.05)$ leaves
\ppnum[3]{confound-w2-l1-tau0p05-unresolved} unresolved, and
$(L,\tau)=(1,0.10)$ leaves
\ppnum[3]{confound-w2-l1-tau0p1-unresolved} unresolved.
The envelope therefore loses informativeness rapidly unless the
characteristic-to-exposure map is tightly constrained, so its classifications
should be read as conditional scenario classifications rather than as
structural systematic-covariance estimates.

% AUTO-GENERATED by pipeline render-paper1-tables -- do not edit by hand.
\begin{table}[H]
\centering
\small
\begin{tabularx}{\linewidth}{rr>{\raggedleft\arraybackslash}X>{\raggedleft\arraybackslash}X>{\raggedleft\arraybackslash}X}
\toprule
$L$ & $\tau_i=\tau_j$ & Robust coverage & Unresolved & Robust violation \\
\midrule
1.00 & 0.00 & 0.707 [0.625, 0.782] & 0.000 [0.000, 0.000] & 0.293 [0.218, 0.375] \\
1.00 & 0.05 & 0.349 [0.262, 0.438] & 0.560 [0.496, 0.627] & 0.091 [0.055, 0.135] \\
1.00 & 0.10 & 0.076 [0.042, 0.115] & 0.898 [0.863, 0.932] & 0.026 [0.012, 0.046] \\
1.25 & 0.00 & 0.015 [0.006, 0.028] & 0.960 [0.941, 0.976] & 0.025 [0.012, 0.045] \\
1.25 & 0.05 & 0.000 [0.000, 0.001] & 0.994 [0.987, 0.998] & 0.006 [0.001, 0.013] \\
1.25 & 0.10 & 0.000 [0.000, 0.000] & 0.999 [0.996, 1.000] & 0.001 [0.000, 0.004] \\
1.50 & 0.00 & 0.000 [0.000, 0.000] & 0.996 [0.991, 0.999] & 0.004 [0.001, 0.009] \\
1.50 & 0.05 & 0.000 [0.000, 0.000] & 0.999 [0.997, 1.000] & 0.001 [0.000, 0.003] \\
1.50 & 0.10 & 0.000 [0.000, 0.000] & 1.000 [1.000, 1.000] & 0.000 [0.000, 0.000] \\
\bottomrule
\end{tabularx}
\caption{Reduced-form fixed-W2 covariance-envelope diagnostic. The article-embedding W2 is the exact balanced empirical transport distance between 128-article measures under the unit-chord ground cost. Return moments are standardized total-return moments, so $c^{\mathrm{actual}}_{ij}=2-2\rho_{ij}$. A dyad is covered when the lower endpoint of the transport-excess interval is nonnegative, violated when its upper endpoint is negative, and unresolved otherwise. Entries are shares of the 4,950 dyads; brackets are percentile intervals from the common multinomial node bootstrap. $L$ and $\tau$ are maintained scales, not estimated from returns; the classifications are not structural systematic-covariance estimates.}
\label{tab:random-exposure-w2}
\end{table}


Before estimating an average association, we ask an intuition-building
question: when two firms are selected because their article distributions are
nearest neighbours within sector, must their subsequent return paths also
track one another?
\Cref{fig:neighbour-outcome-paths} traces cumulative log returns for
\ppnum[0]{neighbour-holdout-n-selected}
within-sector pairs whose firms are Wasserstein nearest neighbours in the
frozen 2018--2020 W$_2$ geometry.
These selected pairs are informative because they make the envelope's
permissive implication visible, but they do not constitute a hypothesis test
or a clean validation sample.

\begin{figure}[htbp]
  \centering
  \ppfigure{0.78}{neighbour_holdout_paths}
  \caption{Illustrative cumulative-log-return paths for
    \ppnum[0]{neighbour-holdout-n-selected} nearest
    within-sector pairs in the frozen 2018--2020 exact W$_2$ matrix.
    Each panel represents one sector, dates run from 4 January 2021 to 30
    December 2022, and solid and dashed lines identify the printed tickers.
    Insets report W$_2$ (characteristic distance), RMS gap (root-mean-square
    path difference), terminal gap (final-date difference), and RMS
    percentile (rank among all within-sector pairs, with 0 denoting the
    closest-tracking pair).
    The selected pairs provide intuition rather than a test, and the
    evaluation window belongs to the broader return panel rather than
    a clean holdout.
  Authors' calculations.}
  \label{fig:neighbour-outcome-paths}
\end{figure}

Some neighbouring pairs track closely over 2021--2022, whereas others
separate materially.
The display therefore does not validate the envelope.
Instead, it clarifies the one-sided implication that textually close firms
\emph{can} comove tightly but need not.
That heterogeneity motivates estimating the average relation over all dyads
rather than drawing an inference from selected pairs.

\subsection{Canonical and timing-separated associations}
\label{sec:post-window-dyadic}

We next ask whether greater article-embedding distance is associated with
weaker return co-movement after symmetric firm effects absorb stable endpoint
heterogeneity, and whether the sign persists when the article-embedding
geometry is held fixed for a later return window.
The canonical estimate answers the first question affirmatively: at
representative dyads across the interquartile range, a pair one standard
deviation farther apart in news content has an implied return correlation
roughly eight points lower.
Return-chord distance maps deterministically to the familiar correlation
scale through \eqref{eq:chord-distance}, because $\hat\rho=1-\hat D^2/2$.
At the median dyad, a chord distance of \ppnum[3]{chord-median} corresponds
to a return correlation of \ppnum[3]{chord-rho-median}.
A one-standard-deviation increase in W$_2$ moves the implied correlation to
\ppnum[3]{chord-rho-after-median}, a change of
\ppnum[3]{chord-delta-rho-median}.
The corresponding changes are \ppnum[3]{chord-delta-rho-q25} at the
lower-quartile dyad and \ppnum[3]{chord-delta-rho-q75} at the
upper-quartile dyad.

\Cref{tab:w2-primary-results} reports the same symmetric-firm-effect
estimator for both return windows.
In the shared 2018--2022 window, the focal coefficient is
\ppnum[3]{confound-w2-coef} return-chord units per unit of W$_2$, with
standard error \ppnum[3]{confound-w2-se} and 95\% interval
[\ppnum[3]{confound-w2-ci-low}, \ppnum[3]{confound-w2-ci-high}].
A one-standard-deviation increase in W$_2$ corresponds to
\ppnum[3]{confound-w2-effect-one-sd} return-chord units, and W$_2$ accounts
for \ppnum[3]{confound-partial-r2} of residual variation relative to
symmetric firm effects.
The positive estimate is consistent with the directional expectation in
\Cref{hyp:monotonicity}.
Endpoint margins are absorbed, but pair-specific common causes remain, so the
coefficient is neither causal nor a numerical estimate of transport excess.

The later-window estimate provides timing-separated corroboration, not
point-in-time out-of-sample validation.
The exercise fixes the 2018--2022 article-embedding geometry before
recomputing return chord distances and firm effects for
\ppvalue{oos-return-panel-date-start} through
\ppvalue{oos-return-panel-date-end}.
The coefficient rises to \ppnum[3]{oos-dyadic-w2-coef}, with 95\% interval
[\ppnum[3]{oos-dyadic-w2-ci-low}, \ppnum[3]{oos-dyadic-w2-ci-high}], and the
standardized effect is \ppnum[3]{oos-dyadic-w2-effect-one-sd}.
The sign persists after separating article selection from the later return
outcomes, but the encoder has no documented training cutoff and common latent
drivers may persist across windows.
Accordingly, the design does not recreate point-in-time feature construction
or establish predictive validation.

% AUTO-GENERATED -- do not edit by hand.
\begin{table}[H]
\centering
\small
\begin{tabularx}{\linewidth}{Xrrrrrrr}
\toprule
Return window & Firms & Dyads & $\hat\beta_{W_2}$ & 95\% CI & $p$ & $\Delta D$/1 SD & Partial $R^2$ \\
\midrule
2018--2022 returns & 100 & 4950 & 1.874 & [1.516, 2.250] & 0.001 & 0.069 & 0.332 \\
2023--2026 returns & 100 & 4950 & 1.945 & [1.516, 2.454] & 0.001 & 0.071 & 0.314 \\
\bottomrule
\end{tabularx}
\caption{Primary symmetric-firm-effect association. The text-side rooted $W_2$ matrix is fixed across rows. The first row relates the 2018--2022 return chord distance to that geometry; the second uses the later 2023--2026 return panel. Intervals and two-sided sign-tail probabilities use the multinomial node bootstrap. $\Delta D$/1 SD standardizes by the full-sample dyadic standard deviation of $W_2$; partial $R^2$ is relative to symmetric additive firm effects. These are conditional associations, not causal estimates or tests of the covariance envelope.}
\label{tab:w2-primary-results}
\end{table}


Having established the average magnitude, we ask whether the canonical
coefficient reflects a visible residual-space relation and whether its sign
depends on functional form or any single firm.
By the Frisch--Waugh--Lovell result, Panel A of
\Cref{fig:dyadic-postfit-diagnostics} residualizes both return distance and
W$_2$ with respect to the symmetric firm effects, so the slope through those
residuals equals the canonical coefficient.
Panels B and C assess functional form and single-firm influence.

\begin{figure}[H]
  \centering
  \ppfigure{1}{dyadic_postfit_diagnostics}
  \caption{Post-fit diagnostics for the canonical
    symmetric-firm-effects specification.
    Panel A plots firm-effect-residualized return chord distance against
    residualized W$_2$; the fitted line uses all dyads, while equal-count bin
    averages are descriptive.
    Panel B plots raw W$_2$ against the residualized outcome and overlays the
    prespecified quadratic sensitivity.
    Panel C reports the ten largest leave-one-firm-out coefficient shifts.
    The diagnostics show a positive residual-space relation without a single
    sign-reversing firm, but they do not remove pair-specific confounding.
  Authors' calculations.}
  \label{fig:dyadic-postfit-diagnostics}
\end{figure}

The diagnostics support the stability of the positive association across
these checks, but neither residualization, the functional-form comparison,
nor leave-one-firm-out analysis identifies the relation or removes
pair-specific confounding.
This distinction motivates asking next whether the association also survives
changes to the representation itself.

\subsection{Representation sensitivity and limits}
\label{sec:representation-sensitivity}

What should readers infer when the article-embedding geometry is compressed,
measured with another encoder or distance metric, or constructed using
encoders trained on a different data vintage?
The sensitivity evidence indicates that the positive association is not
unique to the full-width canonical representation, although sufficiently
severe compression weakens it and the comparisons do not rank representations.
Moderate compression preserves the association, whereas the more severe
compression steps attenuate it, locating a stress boundary at which
discarding semantic coordinates begins to weaken the measured relation.
The Qwen capacity variants and the BGE representation retain positive
associations, and their paired differences from the relevant references are
not distinguishable after adjustment.
W$_1$ nearly preserves the canonical dyad ordering, whereas energy retains a
positive association under a materially different ordering.
The three matched EttaX vintage contrasts likewise meet their declared
equivalence rule within the benchmark-calibrated margin when architecture,
training recipe, compute, tokenizer, sampler, and seed are held fixed across
pre-, within-, and post-sample Wikipedia snapshots.
These results and their paired contrasts are reported in
\Cref{tab:w2-representation-sensitivity}, with construction details and the
article-sampling checks in \Cref{sec:representation-diagnostics}.
The evidence therefore supports descriptive robustness to several
representation choices, not encoder superiority, a causal interpretation, or
the absence of a vintage effect at the much larger capacity of the
production encoder.
These are sensitivity checks rather than additional hypotheses.
